\documentclass[a4paper]{article}


%% Define your personal info here %%%%%%%%%%%%%%%%%%%%%%% <--------------- X
\newcommand\TPid{A}
\newcommand\TPname{Linear Algebra}
\newcommand\Firstname{Donnet}
\newcommand\Familyname{Michel Jean Joseph}
\newcommand\Email{\Firstname.\Familyname@etu.unige.ch}
%%%%%%%%%%%%%%%%%%%%%%%%%%%%%%%%%%%%%%%%%%%%%%%%%%%%%%%

\input{preambule}


%%%%%%%% DOCUMENT %%%%%%%%
\begin{document}
% % % % %
% % % %
% % %
% %
%   

\input{title}

%%% Examples
% #FIGURE

% \begin{figure}[H]
%     \centering
%     \includegraphics[width=0.6\textwidth]{figureName.png}
%     \caption{Caption}
% \end{figure}

% \begin{figure}[H]
%     \centering
%     \begin{subfigure}[t]{0.4\textwidth}
%         \centering
%         \includegraphics[width=\textwidth]{fig_1-1.png}
%         \caption*{1, 1}
%     \end{subfigure}
%     % / \quad / \hfill
%     \begin{subfigure}[t]{0.4\textwidth}
%         \centering
%         \includegraphics[width=\textwidth]{fig_1-2.png}
%         \caption*{1, 2}
%     \end{subfigure}

%     % \hfill \newline

%     \centering
%     \begin{subfigure}[t]{0.25\textwidth}
%         \centering
%         \includegraphics[width=\textwidth]{fig_1-1.png}
%         \caption*{2, 1}
%     \end{subfigure}
%     % / \quad / \hfill
%     \begin{subfigure}[t]{0.25\textwidth}
%         \centering
%         \includegraphics[width=\textwidth]{fig_1-2.png}
%         \caption*{2, 2}
%     \end{subfigure}
%     % / \quad / \hfill
%     \begin{subfigure}[t]{0.25\textwidth}
%         \centering
%         \includegraphics[width=\textwidth]{fig_1-3.png}
%         \caption*{2, 3}
%     \end{subfigure}

%     \caption{Main caption}
% \end{figure}

% \begin{figure}[H]
%     \centering
%     \begin{minipage}[t]{0.4\textwidth}
%         \centering
%         \includegraphics[width=\textwidth]{fig_1.png}
%         \caption*{1}
%     \end{minipage}
%     \begin{minipage}[t]{0.4\textwidth}
%         \centering
%         \includegraphics[width=\textwidth]{fig_2.png}
%         \caption*{2}
%     \end{minipage}
%     \caption{Main caption}
% \end{figure}


% #CODE

% \begin{lstpython}
% MyCode
% \end{lstpython}

% \pythoninline{MyCode}

% #TABLE
% https://www.tablesgenerator.com/#


% #SECTION

% 
% 
%%%%% ================================================================ %%%%%
% \section{Section}

%
%%% =-=-=-=-=-=-=-=-=-=-=-=-=-=-=-=-=-=-=-=-=-=-=-=-=-=-=-=-=-=-=- %%%
% \subsection{Subsection}

%
% ---------------------------------------------------------- %
% \subsubsection{Subsubsection}

% 
% 
%%%%% ================================================================ %%%%%
\section{Theory questions}
Use the slides and your notes to answer these questions. You do not need to write code. Provide short explanations and examples where possible.


\begin{enumerate}
	\item What is a vector? What is a vector space? Give an example of a vector in $\mathbb{R}^2$ or $\mathbb{R}^3$.
    
    \hrulefill\\
    A vector is a mathematical object that have a length, a direction and a sense.

    A vector space is a space of vectors with linear operations such as addition between vectors and multiplication by a scalar.
    It can be generated by the combitation of a few vector, according to the dimension to the space.
    For instance, we can generate all vectors of a vector space of dimension 2 by using 2 vectors.

    \item What do we mean by the ``dimension" $D$ of a dataset?
        
    \hrulefill\\
    The dimension of a dataset indicate the minimum number of vectors useful to generate the entire dataset by combinaisons between this vectors.
    
    \item Create two small datasets: one scalar dataset of 5 numbers, and one vector dataset of 5 vectors in $\mathbb{R}^3$. Compute their mean and variance. Then, write the formal mathematical definitions of mean and variance. 
        
    \hrulefill\\
    Write your answer here

    \item Define the scalar product. Compute it for two vectors from your vector dataset above. How does the scalar product allow us to compute the projection of one vector onto another?
        
    \hrulefill\\
    Write your answer here
    
    \item What is a basis of a vector space? 
        
    \hrulefill\\
    Write your answer here
    
    \item Explain the eigenvalue decomposition (EVD) and singular value decomposition (SVD). Why are they important in data science?
        
    \hrulefill\\
    Write your answer here
    
    \item Give one concrete example of the curse of dimensionality, and one example of the blessing of dimensionality.
        
    \hrulefill\\
    Write your answer here
    
    \item If you divide a $D$-dimensional space into $b$ bins per dimension, and want $k$ samples per bin, how many samples are needed in total?
        
    \hrulefill\\
    Write your answer here
    
    \item Why does data become more sparse as $D$ grows? What happens if the number of samples stays fixed?
        
    \hrulefill\\
    Write your answer here
    
    \item Compare a hypersphere of radius $r$ with its enclosing hypercube $[-r,r]^D$. As $D$ increases, what happens to their volume ratio? What does this imply for uniform sampling?
        
    \hrulefill\\
    Write your answer here
    
    \item In high dimensions, what is the typical angle between two random vectors? How does this angle change as $D$ increases?
        
    \hrulefill\\
    Write your answer here
     
    \item What is meant by a ``concentration phenomenon"? Give one example where it is useful, and one where it is problematic.
        
    \hrulefill\\
    Write your answer here
    
    \item State one concentration inequality (e.g., Markov, Chebyshev, Hoeffding). Explain briefly what it tells us.
        
    \hrulefill\\
    Write your answer here
    
    \item What is a PAC (Probably Approximately Correct) analysis in machine learning?
        
    \hrulefill\\
    Write your answer here

\end{enumerate}

% 
% 
%%%%% ================================================================ %%%%%
\section{Matrix}

\begin{enumerate}
	\item Find a quadratic polynomial, say $f(x) = ax^2 + bx + c$, such that
	      \[
		      f(1) = -5, \quad f(2) = 1, \quad f(3) = 11.
	      \]
	      \hrulefill\\
	      Write your solution here


	\item Let $a$, $b$ be some fixed parameters. Solve the system of linear equations
	      \[
		      \begin{cases*}
			      x + ay = 2 \\
			      bx + 2y = 3
		      \end{cases*}
	      \]
	      \hrulefill\\
	      Write your solution here

	\item Find the inverse of the following matrix, if it exists:
	      $$
		      A = \begin{pmatrix}
			      1 & 2 & 3 \\
			      0 & 1 & 4 \\
			      5 & 6 & 0
		      \end{pmatrix}
	      $$
	      \hrulefill\\
	      Write your solution here

	\item Solve the following matrix equation for \( X \):
	      $$
		      AX = B, \quad \text{where} \quad A = \begin{pmatrix} 1 & 3 \\ 2 & 5 \end{pmatrix}, \quad B = \begin{pmatrix} 7 & 8 \\ 9 & 10 \end{pmatrix}.
	      $$
	      \hrulefill\\
	      Write your solution here

\end{enumerate}

% 
% 
%%%%% ================================================================ %%%%%
\clearpage
\section{The importance of the mathematical concept behind a code}

\begin{enumerate}
	\item Download, open and run several times the PYTHON file "some\_script.py".

	      Explain the function \mintinline{python}{def project_on_first(u, v)} in geometrical terms and then explain the results of this code with clear mathematical concepts (you have to talk about projections and scalar product). \emph{It is very important to understand the mathematical concepts behind a program!!}

	      \hrulefill\\
	      Write your solution here

	\item Consider $u = (u_i)_{1 \le i \le 7}$ and $v = (v_i)_{1 \le i \le 7}$, two vectors of length 7. How would you clearly write the following part of the script using some mathematical notation?

	      \begin{minted}
    [frame=lines,
    		framesep=2mm,
    		linenos]
    {python}
    # import numpy library as np in order to work on tensors
    import numpy as np
    # create 2 random vectors of length 7
    u = np.random.randn(7)
    v = np.random.randn(7)
    # perform some computations
    r = 0
    for ui, vi in zip(u, v) :
        r += ui * vi
    \end{minted}
	      Name the mathematics behind and rewrite the last 3 lines of code in one.

	      \hrulefill\\
	      Write your solution here


	\item Complete the code to construct a vector $v$ orthogonal to the vector $u$ and with the same norm as $u$. 
	      Comment each line of your code.

	      \hrulefill\\
	      Write your solution here\

       \item Given two vectors \( u \) and \( v \) in \( \mathbb{R}^n \), write Python code to compute the cosine of the angle between them. Explain the mathematical concept behind this computation, and show how you would write it in one equation.

       \hrulefill\\
	      Write your solution here\

\end{enumerate}

% 
% 
%%%%% ================================================================ %%%%%
\clearpage
\section{Computing Eigenvalues, Eigenvectors, and Determinants}

\begin{enumerate}
	\item Find the determinant, eigenvectors, and eigenvalues of the matrix
	      \[
		      \begin{bmatrix}
			      5  & 6 & 3  \\
			      -1 & 0 & 1  \\
			      1  & 2 & -1
		      \end{bmatrix}
	      \]
	      Compare your results with those of the computer.

	      \hrulefill \\
	      Write your solution here


	\item The covariance matrix for n samples $x_1, \dots, x_n$ , each represented by a $d \times 1$ column vector, is given by
	      \[
		      C = \frac{1}{n - 1} \sum^n_{i=1} (x_i - \mu)(x_i - \mu)^T,
	      \]
	      where $C$ is a $d \times d$ matrix and $\mu = \sum^n_{i=1} x_i/n$ is the \textit{sample mean}.
	      Prove that $C$ is always positive semidefinite. (Note: A symmetric matrix $C$ of size $d \times d$ is \textit{positive semidefinite} if $v^T C v \geq 0$ for every $d \times 1$ vector $v$.)

	      \hrulefill \\
	      Write your solution here

	\item In this portion of the exercise, we will calculate the eigenvalues of the covariance matrices of six data sets listed as follows:
	      \begin{table}[H]
		      \centering
		      \begin{tabularx}{.7\textwidth}{rccX}
			      \hline
			      filename                          & $n$  & $d$ & description                                                          \\ \hline
			      \texttt{tp1\_artificialdata[1-3]} & 1024 & 100 & Artificial data generated from various auto-regression (AR-1) models \\
			      \texttt{tp1\_artificialdata4}     & 1024 & 100 & Random Gaussian data                                                 \\
			      \texttt{tp1\_freyfaces}           & 1965 & 560 & Facial images of a man named Brendan                                 \\
			      \texttt{tp1\_digit2}              & 5958 & 784 & Hand-written images of "2"                                           \\ \hline
		      \end{tabularx}
	      \end{table}
	      To access each data set, download and unzip \texttt{data for python.zip}.
	      Each file contains a $n \times d$ data matrix with rows representing $n$ different samples in $\mathbb{R}^d$
	      For example, \texttt{tp1\_artificialdata1} contains a data matrix of size $1024 \times 100$ (1024 samples, 100 features).
	      Once each dataset has been imported into \texttt{Python}, complete the following tasks:

	      \begin{enumerate}[label=(\alph*)]
		      \item Compute the covariance matrix of each dataset;

		      \item Compute the eigenvalues for each covariance matrix;

		      \item Compute the determinant of each covariance matrix;

		      \item Compute the product of the eigenvalues for each covariance matrix;

		      \item Display the eigenspectrum for each covariance matrix in a 2D plot, where the $x$-axis shows the rank of the eigenvalues, ranging from 1 (the largest eigenvalue) to 100 (the 100-th largest eigenvalue), and the $y$-axis shows the corresponding eigenvalue.
	      \end{enumerate}
	      Describe the relationship between the product of the eigenvalues and the
determinant of each covariance matrix. In addition, describe the observations
you made regarding the plot of the spectrum of eigenvalues of each covariance
matrix. How can you explain the disparities between datasets?

	      \hrulefill \\
	      Write your solution here

       \item Let \( A \in \mathbb{R}^{n \times n} \) be a square matrix. Show that if \( A \) has an eigenvalue \( \lambda = 0 \), then the matrix \( A \)  is singular. Furthermore, prove that the determinant of \( A \) must be zero if any eigenvalue is zero.

       \hrulefill \\
	      Write your solution here
\end{enumerate}

%%%%% ================================================================ %%%%%
\clearpage
\section{Computing Projection Onto a Line}
There are given a line $\alpha : 3x + 4y = -6$ and a point $A$ with the coordinates (-1, 3).
\begin{enumerate}
	\item Find a distance from the point $A$ to the line $\alpha$ using linear algebra (not coordinate method).
	\item Explain your code using a sketch and some mathematics (there should be a scalar product somewhere).
\end{enumerate}

\hrulefill \\
Write your solution here


% 
% 
%%%%% ================================================================ %%%%%
% Note, the assessment is mainly based on your report, which should include your answers to all questions and the experimental results.
% \textit{Importance is given on the mathematical explanations of your works and your codes should be commented}
\end{document}

